\section{Résumé exécutif}

Cette page répond à quatre questions sans jargon : que fait le système, où en
est-il, qu'est-ce qui est prouvé, et que faut-il décider.

\subsection{Ce que le système fait}

LANCE conduit un audit de sécurité sur un réseau d'objets connectés \emph{dont
la propriété et l'autorisation sont établies}. Il explore le réseau, émet des
hypothèses de failles, tente de les vérifier par des preuves observables, puis
rédige un rapport. Deux politiques peuvent mener cet audit : une automatisation
à règles et un agent à base de modèle de langage. Le projet cherche à savoir
laquelle des deux réussit mieux, à outils, preuves et budgets identiques.

\subsection{Où en est le projet}

\begin{table}[H]
\centering
\caption{État en une ligne par volet.}
\label{tab:etat-volet}
\small
\begin{tabular}{@{}p{3.3cm}p{1.9cm}p{8.3cm}@{}}
\toprule
Volet & État & Précision \\
\midrule
Moteur d'audit & mature & Six phases en service, quatre modes, frontière d'exécution qui journalise chaque appel \\
Preuves & mature & Contrat briefing/oracle/évaluateur appliqué à cinq familles ; le journal fait foi, pas la décision \\
Harnais & consolidé & Noyau partagé depuis le 29 octobre ; bac à sable exécutable hors laboratoire \\
Infrastructure & rétablie & Trois instances, circuit de livraison automatique, sept défauts corrigés \\
Benchmark & partiel & Une hypothèse mesurée sur six, une famille rejetée, aucune inférence calculable \\
Preuve scientifique & absente & Aucune supériorité du modèle n'est établie à ce jour \\
\bottomrule
\end{tabular}
\end{table}

\subsection{Ce qui est prouvé, et ce qui ne l'est pas}

\textbf{Prouvé.} Chaque appel d'outil est tracé et borné ; une preuve perdue
interrompt le run. Un harnais peut tourner hors du laboratoire, sous un garde-fou
dont les angles morts sont déclarés. Sur la famille la plus documentée, un
troisième script --- qui lit par convention, sans modèle --- obtient le même
résultat que l'agent, ce qui \emph{absorbe} le gain attribué au modèle. Le corpus
est majoritairement détectable par signature, ce qui rend un écart nul attendu
par construction.

\textbf{Non prouvé.} Qu'un modèle fasse mieux qu'un script compétent. Qu'un
écart soit significatif. Que les résultats se généralisent au-delà de huit
dossiers et d'une propriété.

\subsection{Les trois décisions qui commandent la suite}

\begin{enumerate}
\item \textbf{L'outil d'inférence.} Il produit aujourd'hui des p-values qui ne
distinguent pas quatre séries de comportements différents. À corriger pour qu'il
refuse de produire un chiffre, ou à retirer du chemin de décision.
\item \textbf{La condition d'ouverture.} Plusieurs familles restent indécidables
faute d'une automatisation de référence écrite. L'écrire, ou renoncer
explicitement à ces familles.
\item \textbf{Les familles indépendantes.} En produire trois, ou publier un
résultat méthodologique borné en le déclarant tel. C'est cette décision qui
détermine si une inférence pourra exister un jour.
\end{enumerate}

\subsection{Comment lire ce rapport}

La partie~\ref{sec:architecture-pipeline} décrit le chemin d'un audit.
La partie~\ref{sec:protocoles} décrit ce qui rend un résultat vérifiable.
La partie~\ref{sec:harnais-sandbox} décrit l'outillage expérimental.
La partie~\ref{sec:infrastructure} décrit la machine et la livraison.
La partie~\ref{sec:benchmark} établit ce qui est mesuré.
La conclusion confronte ces volets et relève leurs contradictions.

Chaque affirmation renvoie au fichier qui la porte. Les comportements que le code
ne permet pas d'établir sont signalés comme tels plutôt que comblés.

\clearpage
\section{Vue d'ensemble du système}
\label{sec:vue-ensemble}

Avant d'entrer dans les contrats, une carte du système et du laboratoire. Elle
situe les composants décrits ensuite et montre où passent les preuves.

\begin{figure}[H]
\centering
\begin{tikzpicture}[
  node distance=8mm and 9mm,
  box/.style={draw=line, fill=white, rounded corners=2pt, line width=.55pt,
    text=ink, align=center, inner sep=5pt, font=\sffamily\footnotesize,
    text width=29mm, minimum height=12mm},
  focal/.style={box, draw=accent, fill=rose, line width=.75pt},
  flow/.style={-{Latex[length=1.8mm,width=1.1mm]}, draw=muted, line width=.6pt}
]
\node[box] (lab) {Laboratoire IoT\\simulé};
\node[box, right=of lab] (run) {Moteur d'audit\\six phases};
\node[focal, right=of run] (tools) {Frontière d'exécution\\journal des preuves};
\node[box, below=of lab] (eval) {Évaluateur\\hors ligne};
\node[box, right=of eval] (verdict) {Verdict};
\node[box, above=of eval] (art) {Artefacts\\du run, 01 à 06};
\node[box, above=of verdict] (oracle) {Oracle privé\\fixé avant le run};

\draw[flow] (lab) -- (run);
\draw[flow] (run) -- (tools);
\draw[flow] (run) -- (art);
\draw[flow] (art) -- (eval);
\draw[flow] (oracle) -- (verdict);
\draw[flow] (eval) -- (verdict);
\draw[flow, dashed] (tools.south) -- (verdict.north east);
\end{tikzpicture}
\caption{Le trajet d'une preuve. Les outils produisent des observations
journalisées ; l'oracle privé, écrit avant l'exécution, permet de recomputer le
verdict depuis le journal sans croire la décision du système.}
\label{fig:tripwire}
\end{figure}

Trois propriétés se lisent sur cette carte.

\textbf{Le journal est la source.} Un résultat n'est pas la conclusion écrite par
la politique, mais ce que l'évaluateur recalcule depuis les observations
archivées. Une preuve citée qui ne figure pas dans le journal est rejetée.

\textbf{L'oracle précède le run.} Il est fixé avant que la politique ne tourne,
et un constat qui le contredit invalide l'essai au lieu de le corriger. Cette
antériorité est ce qui empêche d'ajuster la réponse après coup.

\textbf{Le bac à sable n'est pas la mesure.} Il exécute les bras déterministes
sans modèle ni laboratoire ; il instrumente le contraste, il ne le mesure pas.

\begin{figure}[H]
\centering
\begin{tikzpicture}[
  node distance=3mm and 5mm,
  h/.style={draw=line, fill=white, rounded corners=2pt, line width=.55pt,
    text=ink, align=center, inner sep=3pt, font=\sffamily\footnotesize},
  dev/.style={h, fill=rose, draw=accent},
  flow/.style={-{Latex[length=1.6mm,width=1mm]}, draw=muted, line width=.6pt}
]
\node[h, text width=22mm] (poste) {Poste de\\développement};
\node[h, text width=24mm, right=of poste] (ci) {GitHub Actions\\validation};
\node[dev, text width=26mm, right=of ci] (lxc) {Conteneur maître\\trois instances};
\node[h, text width=22mm, right=of lxc] (umons) {Relais UMONS\\modèle};

\draw[flow] (poste) -- node[above, font=\sffamily\tiny, text=muted] {push} (ci);
\draw[flow] (ci) -- node[above, font=\sffamily\tiny, text=muted] {déploie} (lxc);
\draw[flow] (lxc) -- node[above, font=\sffamily\tiny, text=muted] {appelle} (umons);
\end{tikzpicture}
\caption{Le circuit de livraison. Aucun travail d'exécution n'a lieu sur le poste
de développement : le push déclenche la validation, qui déploie sur le conteneur.}
\label{fig:circuit}
\end{figure}

\clearpage
